% ===================================================================
%  કોષ્ટક નં. ૬ — ઘરની અંદર, રાચરચીલું, ખોરાક, પ્રકાશ.
%
%  Traduction, faite en 2026, en gujarati d'aujourd'hui, sur l'ido de
%  texto/io. Regles de la colonne : voir 10-kostak-01.tex. Cinq
%  scenes, cinq numerotations. C'est le tableau le plus charge du
%  livret : 211 renvois, et quarante-sept paires annoncees par
%  parigi.py.
%
%  LA CUISINE EST GUJARATIE ET LE SALON NE L'EST PAS, ET C'EST LE
%  MEME PARTAGE QUE DANS LA COLONNE HAOUSSA AU MEME TABLEAU :
%
%    ચૂલો (1)      le foyer        તપેલી (43)    la marmite
%    ડોયો (21)     la louche       ઝારો (22)     l'ecumoire
%    કડાઈ (47)     la poele        ચાળણી (28)    la passoire
%    છીણી (29)     la rape         ગળણી (38)     l'entonnoir
%    સાવરણી (34)   le balai        ભઠ્ઠી (20)    le four
%    દીવાસળી (12)  les allumettes  મીઠું (14)     le sel
%
%  Contre, au salon : સોફો, આલ્બમ, પિયાનો, ફાનસ, લામ્બ્રકિન. Pas un
%  objet de cette piece-la n'existait ici avant qu'on l'y porte.
%
%  ET LE CHARBON DE TERRE (17) EST « પથ્થરનો કોલસો », LE CHARBON DE
%  PIERRE, CE QUI EST EXACTEMENT LE MOT HAOUSSA. « gawayin dutse »,
%  le charbon de bois de pierre, avait ete releve au meme numero du
%  meme tableau, et les deux langues n'ont jamais eu le moindre
%  contact. Elles avaient toutes les deux le charbon de bois et ont
%  nomme l'autre par la matiere. Ce n'est pas un emprunt de l'une a
%  l'autre : c'est le meme raisonnement fait deux fois.
%
%  MAIS LE TAPIS (40) EST « ગાલીચો », ET C'EST LE MOT PERSAN
%  قالیچه, LETTRE POUR LETTRE. C'est le seul objet du livret sur
%  lequel deux colonnes du releve disent litteralement le meme mot
%  sans que l'une l'ait pris a l'autre par hasard : le tapis EST
%  persan, l'objet et le nom sont partis ensemble de Kirman vers
%  l'Inde, et le gujarati n'a jamais eu besoin d'en fabriquer un
%  autre.
%
%  TROISIEME REFUS DE LA COLONNE : LE LIT (17). Le gujarati a deux
%  mots. « ખાટલો » est le lit de corde tressee sur quatre pieds de
%  bois, qu'on sort dans la cour le soir et qu'on rentre le matin, et
%  c'est le lit de la maison gujaratie. « પલંગ » est persan et
%  designe le lit a chassis, monte, garni. La planche grave un
%  chassis, un sommier a ressorts, un matelas de laine, un traversin,
%  un edredon et un ciel de lit : on ecrit પલંગ. Une langue qui a
%  deux mots pour un meuble n'a pas a choisir le plus sien, elle doit
%  choisir le bon.
%
%  QUATRIEME REFUS : LE PAIN (12). « રોટલી » est la galette de ble
%  qu'on roule au વેલણ et qu'on cuit sur la તવી, deux fois par jour,
%  dans toutes les maisons. Le boulanger de la planche a un four et
%  en tire une miche. On ecrit « બ્રેડ ». C'EST LE MEME REFUS QUE LA
%  COLONNE HAOUSSA A FAIT AU TABLEAU 15 avec « gurasa », le pain de
%  froment cuit sur une plaque : deux langues, deux galettes, un seul
%  four europeen, et le meme mot etranger ecrit a la fin.
%
%  ET LE CHAUFFE-BAIN (2) EST « બંબો », DU PORTUGAIS bomba —
%  CINQUIEME MOT PORTUGAIS DE LA COLONNE, apres મેજ la table, પાદરી
%  le pretre, અનનાસ l'ananas et મિસ્ત્રી l'ouvrier qualifie. Cinq
%  tableaux, cinq mots, et pas un qui soit venu par le meme chemin :
%  un meuble, un pretre, un fruit d'Amerique, un metier et une
%  machine.
%
%  QUATRE FAUTES D'ORDRE SUR QUARANTE-SEPT PAIRES ANNONCEES, ET LES
%  QUATRE SONT DE LA MEME FORME : un participe ou une preposition de
%  l'ido posait le complement AVANT son support — « la laborkorbo
%  proxim la tabureto », « pikturi akrochita an la parieto », « tasi
%  e subtasi sur la plad-moblo », « kaldrono pozita sur la
%  karel-sulo ». Chaque fois le gujarati avait ramene le support
%  devant, ce qu'il fait naturellement. Chaque fois la relative en
%  « જે » a suffi a le rejeter derriere. LE REMEDE EST TOUJOURS LE
%  MEME DEPUIS LE TABLEAU 3, ET C'EST LA SEULE CHOSE QU'IL FAILLE
%  SAVOIR POUR ECRIRE CETTE COLONNE.
%
%  ENFIN LE THE (31) EST « ચા » ET LE CAFE (63, 64) EST « કૉફી », ET
%  LES DEUX MOTS PORTENT LEUR ROUTE. ચા vient du chinois chá par la
%  terre, a travers la Perse — c'est le meme cha que le persan چای et
%  que le russe чай. કૉફી vient de l'arabe qahwa par la MER, avec les
%  Anglais. Deux boissons, deux routes, et une langue qui a garde
%  dans chaque mot celle par laquelle il est arrive.
% ===================================================================

%%K t06-apar-1 apar
\VUsaut{6.0mm}
\VUtitre{15.0pt}{60}{કોષ્ટક નં. ૬}
\VUsaut{1.4mm}
\VUtitre{11.4pt}{0}{ઘરની અંદર, રાચરચીલું, ખોરાક,}
\VUsaut{0.4mm}
\VUtitre{11.4pt}{0}{પ્રકાશ.}
\VUsaut{2.2mm}

%%K t06-apar-2 apar
ઘર પૂરું થયું છે. યોહાનેસના કાકા પોતાના નવા ઘરમાં રહેવા આવ્યા છે,
જેની ગોઠવણ આપણે જાણીએ છીએ. હવે જોઈએ કે તેમણે પોતાનું ઘર કેવી રીતે
સજાવ્યું છે. પહેલાં આપણે તપાસીશું :

%%K t06-tit-1 sub
\VUpk{10.2pt}{40}{શયનખંડ.}
\VUsaut{3.0mm}

%%K t06-c1-01-1 p
1. --- આપણે અયોગ્ય વખતે આવીએ છીએ, કેમ કે ઓરડાની નોકરાણી
\VUgras{ઓરડો} સાફ કરવામાં રોકાયેલી છે.

%%K t06-c1-02-1 p
2. --- \VUgras{શણગારમેજ} \textsuperscript{(1)} પર અહીંતહીં જુદી જુદી
વસ્તુઓ છે જેનાથી માણસ ધુએ છે, વાળ ઓળે છે, ટૂંકમાં શણગાર કરે છે. તે છે
\VUgras{કૂંડી} \textsuperscript{(2)} અને \VUgras{પાણીનો ઘડો}
\textsuperscript{(3)}, \VUgras{વાળનો બ્રશ} \textsuperscript{(4)},
\VUgras{કાંસકો} \textsuperscript{(5)}, \VUgras{સાબુ}
\textsuperscript{(6)} અને \VUgras{સ્પોન્જ} \textsuperscript{(7)}, છેલ્લે
દાંતના પ્રવાહી મંજન માટે \VUgras{નાની શીશી} \textsuperscript{(8)}.

%%K t06-c1-03-1 p
3. --- જમીન પર, શણગારમેજ આગળ, આ રહી ગંદું પાણી ભેગું કરવા માટેની
\VUgras{ડોલ} \textsuperscript{(9)}.

%%K t06-c1-04-1 p
4. --- \VUgras{આગલા ઓરડામાં} \textsuperscript{(10)} આપણે કેટલીક
\VUgras{કપડાંની ખીંટીઓ} \textsuperscript{(13)} અને બે \VUgras{છત્રી}
\textsuperscript{(11)} જોઈએ છીએ \VUgras{છત્રીદાનીમાં}
\textsuperscript{(12)}.

%%K t06-c1-05-1 p
5. --- બારણાની બીજી બાજુએ \VUgras{અરીસાવાળું કબાટ}
\textsuperscript{(14)} છે, જેમાં \VUgras{કપડાં} \textsuperscript{(15)}
છે.

%%K t06-c1-06-1 p
6. --- \VUgras{ઓરડાની નોકરાણી} \textsuperscript{(6)} \VUgras{પલંગ}
\textsuperscript{(17)} ગોઠવે છે તેનો લાભ લઈને આપણે તેના જુદા જુદા ભાગ
તપાસીએ. \VUgras{પલંગની ચોકઠીમાં} \textsuperscript{(20)} સારી
\VUgras{સ્પ્રિંગની ગાદી} \textsuperscript{(21)} છે, જેના પર યુસ્તિના
હમણાં \VUgras{ઊનની ગાદી} \textsuperscript{(22)} મૂકશે. પછી તે
ખુરશીઓ પરથી બે \VUgras{ચાદરો} \textsuperscript{(23)} અને
\VUgras{ધાબળો} \textsuperscript{(24)} લેશે. પછી તે પલંગના માથા પર
\VUgras{લાંબું ઓશીકું} \textsuperscript{(25)} અને \VUgras{ઓશીકું}
\textsuperscript{(26)} મૂકશે, જે \VUgras{રજાઈ} \textsuperscript{(27)}
સાથે \VUgras{પલંગની જાજમ} \textsuperscript{(32)} પર છે. તે
\VUgras{પડદા} \textsuperscript{(19)} અને \VUgras{પલંગનું છત્ર}
\textsuperscript{(18)} ની ધૂળ ખંખેરવા પીંછાંનો ઝૂડો વાપરે છે, અને આ
રહ્યો કામનો એક ભાગ પૂરો.

%%K t06-c1-07-1 p
7. --- પછી તેણે \VUgras{પલંગ પાસેનું નાનું મેજ} \textsuperscript{(28)}
થોડું ગોઠવવું પડશે, \VUgras{જગાડનારું ઘડિયાળ} \textsuperscript{(29)}
ચાવી દેવી પડશે, \VUgras{રાતનો નાનો દીવો} \textsuperscript{(30)} તૈયાર
કરવો પડશે, \VUgras{મીણબત્તી} \textsuperscript{(31)} કાઢીને દાની સાફ
કરવી પડશે.

%%K t06-tit-2 sub
\VUsaut{5.0mm}
\VUpk{10.2pt}{40}{કામની ટોપલી અને સગડીનો સામાન.}
\VUsaut{3.0mm}

%%K t06-c1-08-1 p
8. --- \VUgras{કામની ટોપલી} \textsuperscript{(34)}, જે \VUgras{ટેભા}
\textsuperscript{(33)} પાસે છે, તેને પણ ગોઠવવાની ઘણી જરૂર છે. તેણે
\VUgras{ભરતકામ} \textsuperscript{(35)} વાળવું પડશે, \VUgras{કામના
મેજમાં} \textsuperscript{(38)} \VUgras{આંગળીનું ટોપકું,}
\VUgras{દોરો} \textsuperscript{(36)}, \VUgras{સોય} \textsuperscript{(37)}
અને \VUgras{કાતર} \textsuperscript{(39)} મૂકવાં પડશે અને મેજનું ઢાંકણ
\VUgras{ચાવીથી} \textsuperscript{(40)} બંધ કરવું પડશે.

%%K t06-c1-09-1 p
9. --- વચ્ચેના \VUgras{એક પગવાળા મેજ} \textsuperscript{(41)} પર
યુસ્તિના \VUgras{પાણીની બાટલીનું} \textsuperscript{(42)} પાણી બદલશે.
તે \VUgras{ખાંડ} મૂકશે \VUgras{ખાંડદાનીમાં} \textsuperscript{(43)},
\VUgras{ખાંડની ચીપિયો} \textsuperscript{(44)} સાફ કરશે અને
\VUgras{કપડાંનો બ્રશ} \textsuperscript{(46)} લઈ જશે.

%%K t06-c1-10-1 p
10. --- ઓરડાની જમણી બાજુ તેની પાસે ઓછું કામ માગશે, કેમ કે
\VUgras{લાંબી આરામખુરશી} \textsuperscript{(47)} અને તેનું
\VUgras{ઓશીકું} \textsuperscript{(48)} પોતાની બરાબર જગ્યાએ છે, અને, ઠંડી
ઋતુ ન હોવાથી, \VUgras{સગડીના} \textsuperscript{(49)} સામાનમાં કશું
ખસેડ્યું નથી. \VUgras{પાવડો} \textsuperscript{(50)}, \VUgras{મોટા
ચીપિયા} \textsuperscript{(51)}, \VUgras{આગના દાંડા}
\textsuperscript{(55)}, \VUgras{રાખની જાળી} \textsuperscript{(56)} સ્વચ્છ
છે ; \VUgras{આગની જાળી} \textsuperscript{(54)} ખસેડી નથી અને
\VUgras{લાકડાની પેટી} \textsuperscript{(52)} \VUgras{બળતણનાં લાકડાંથી}
\textsuperscript{(53)} સારી રીતે ભરેલી છે.

%%K t06-noto-1 noto
\VUnotes{143.2mm}{%
(*) Babo = નાનું બાળક, E. baby, F. bébé, I. bambino.%
}

%%K t06-noto-2 noto
\VUnotes{143.2mm}{%
(*) Lambrequino = F. lambrequin, S. lambrequines, Port. lambrequins,
D. E. પણ lambrequins, એટલે D. E. F. P. S.%
}

%%K t06-c1-11-1 p
11. --- તેમ છતાં તેણે \VUgras{લોલકવાળા ઘડિયાળની}
\textsuperscript{(57)}, \VUgras{મીણબત્તીદાનીઓની} \textsuperscript{(58)}
અને \VUgras{શણગારની વસ્તુઓની} \textsuperscript{(59)} ધૂળ ખંખેરવી પડશે,
છેલ્લે \VUgras{અરીસો} \textsuperscript{(60)} લૂછવો પડશે.

%%K t06-c1-12-1 p
12. --- નાના બાળકના (બાબાના (*)) \VUgras{પારણા} \textsuperscript{(61)}
વિશે તો, તે તૈયાર જ છે.

%%K t06-tit-3 sub
\VUsaut{5.0mm}
\VUpk{10.2pt}{40}{બેઠકખંડ.}
\VUsaut{3.0mm}

%%K t06-c2-01-1 p
1. --- હવે આ રહ્યો બેઠકખંડ. તે ઘણો સુંદર \VUgras{ઓરડો} છે. જમણા
ખૂણામાં આવેલી સગડીને નાજુક \VUgras{નાની પ્રતિમા} \textsuperscript{(1)}
અને બે સુંદર \VUgras{ફૂલદાની} \textsuperscript{(3)} શણગારે છે, અને
મખમલના \VUgras{લામ્બ્રકિનથી} \textsuperscript{(2)} (*) ઢાંકી છે.

%%K t06-c2-02-1 p
2. --- સગડીના તણખા ગાલીચાને બાળે નહીં તે માટે સગડી આગળ તાંબાની
\VUgras{તણખા-જાળી} \textsuperscript{(4)} મૂકી છે.

%%K t06-c2-03-1 p
3. --- પાસે, સ્ત્રીઓનું સુઘડ \VUgras{લેખનમેજ} \textsuperscript{(5)}
ખુલ્લું છે. તેના પર \VUgras{ઘરની શેઠાણી} \textsuperscript{(23)}
પોતાના પત્રો લખે છે.

%%K t06-c2-04-1 p
4. --- બારીને \VUgras{બારીના પડદા} \textsuperscript{(6)} શણગારે છે,
તેમની \VUgras{દોરીઓ} \textsuperscript{(8)} સાથે, જે \VUgras{હૂક}
\textsuperscript{(7)} પર ભરાવેલી છે, અને કાપડના મોટા પડદા \VUgras{ઉપર
ભેગા કરેલા} \textsuperscript{(9)}.

%%K t06-c2-05-1 p
5. --- બારીના ગોખલામાં, \VUgras{કૂંડાના ઠામમાં}
\textsuperscript{(11)} લીલો \VUgras{છોડ} \textsuperscript{(10)} છે,
સુંદર તાડ, જેનાં પાંદડાં લગભગ \VUgras{ઘાસલેટના દીવા}
\textsuperscript{(12)} સુધી ફેલાય છે.

%%K t06-c2-06-1 p
6. --- દીવાનું \VUgras{છત્ર} \textsuperscript{(13)} મારિયાએ ગોઠવ્યું છે,
મોહક \VUgras{જુવાન છોકરી} \textsuperscript{(14)}, જે પિયાનો
\textsuperscript{(15)} પાસે છે. \VUgras{ટેભા} \textsuperscript{(16)} પર
ટટ્ટાર બેસીને, તે સંગીતનો અભ્યાસ કરે છે. તે ઘણી કુશળ અને કાળજીવાળી છે,
જેમ તેનું \VUgras{સંગીતનું કબાટ} \textsuperscript{(17)} સાબિત કરે છે, જે
સંપૂર્ણ રીતે ગોઠવેલું છે.

%%K t06-tit-4 sub
\VUsaut{5.0mm}
\VUpk{10.2pt}{40}{તોફાની છોકરો.}
\VUsaut{3.0mm}

%%K t06-c2-07-1 p
7. --- તેનો ભાઈ, પાઉલુસ, \VUgras{પુસ્તક} \textsuperscript{(19)} શોધે છે
\VUgras{પુસ્તકના કબાટમાં} \textsuperscript{(18)}. તે \VUgras{જુવાન}
\textsuperscript{(20)} છે, ચંચળ અને અસહ્ય. જે પુસ્તક બહુ ઊંચે છે તેને
પહોંચવા માટે કબાટને વળગીને, તે ઉપર રહેલી \VUgras{પ્રતિમા}
\textsuperscript{(21)} પાડવાનું જોખમ વહોરે છે.

%%K t06-c2-08-1 p
8. --- આ મૂર્ખાઈ કરવા માટે તે એ વાતનો લાભ લે છે કે તેની મા એક સખી
સાથે વાત કરવામાં રોકાયેલી છે, એક \VUgras{સ્ત્રી} \textsuperscript{(22)}
જે તેમને ત્યાં થોડા દિવસ ગાળવા આવી છે. મા \VUgras{સોફા}
\textsuperscript{(24)} પર બેઠી છે \VUgras{આરામખુરશી}
\textsuperscript{(25)} પાસે, અને તેની સખી \VUgras{બેઠક}
\textsuperscript{(27)} પર છે, જેનો માત્ર \VUgras{પીઠનો ટેકો}
\textsuperscript{(26)} જ દેખાય છે.

%%K t06-c2-09-1 p
9. --- ભીંતે લટકતા મોટા \VUgras{ચોકઠામાં} \textsuperscript{(30)} એક
સગીનું \VUgras{ચિત્ર} \textsuperscript{(29)} છે. મેજ પર મૂકેલા
\VUgras{આલ્બમમાં} \textsuperscript{(32)} કુટુંબનાં બાકીનાં માણસોના
\VUgras{ફોટા} \textsuperscript{(33)} ભેગા કર્યા છે. છોકરાંઓએ હમણાં જ તે
ફેરવ્યું છે, કેમ કે \VUgras{ખુરશીઓ} \textsuperscript{(31)} હજી મેજની
આસપાસ ભેગી છે.

%%K t06-c2-10-1 p
10. --- ચિનાઈ માટીની \VUgras{ચીની ફૂલદાની} \textsuperscript{(34)}, જે
\VUgras{કાગળ કાપવાની છરી} \textsuperscript{(35)} પાસે છે, મારિયાને તેના
તહેવાર માટે ભેટ મળી હતી, અને ઓરડાના ખૂણાને શણગારતો \VUgras{પડદો}
\textsuperscript{(36)} તેના કાકા ચીનથી લાવ્યા હતા.

%%K t06-c2-11-1 p
11. --- ઘણાં સુંદર \VUgras{ચિત્રોથી} \textsuperscript{(37)} આનંદિત,
જે \VUgras{ભીંતે} \textsuperscript{(42)} લટકે છે, \VUgras{છતના
ઝુમ્મરથી}
\textsuperscript{(38)} પ્રકાશિત, જેના \VUgras{વીજળીના દીવા}
\textsuperscript{(39)} સાંજે આનંદથી ચમકે છે, આ બેઠકખંડ ઘણો સુખદ લાગે
છે. લાકડાના ભોંયતળિયા પર પાથરેલો નરમ \VUgras{ગાલીચો}
\textsuperscript{(40)}, અને એટલી સહેલી \VUgras{વીજળીની ઘંટડી}
\textsuperscript{(41)}, તેની સગવડ પૂરી કરે છે.

%%K t06-tit-5 sub
\VUsaut{3.6mm}
\VUpk{10.2pt}{40}{ભોજનખંડ.}
\VUsaut{3.0mm}

%%K t06-c3-01-1 p
1. --- રાતના ભોજનનો ઘંટ વાગ્યો. મારિયા સિવાય આખું કુટુંબ મેજની આસપાસ
ભેગું થયું છે. ભોજનખંડ ચિનાઈ માટીની ઉત્તમ \VUgras{સગડીથી}
\textsuperscript{(1)} સુખદ રીતે ગરમ છે, જેને પાતળી \VUgras{નળી}
\textsuperscript{(2)} છે.

%%K t06-c3-02-1 p
2. --- આ ઓરડામાં બહુ રાચરચીલું નથી. ભીંતે, \VUgras{ટેભા}
\textsuperscript{(4)} ની ઉપર, શણગારનો \VUgras{નાનો ફુવારો}
\textsuperscript{(3)} લટકે છે. ઘણી પાસે, \VUgras{વાસણનું કબાટ}
\textsuperscript{(5)} છે જેના પર ઠંડી વાનગીઓ, મીઠાઈની થાળીઓ, અને મેજનાં
જુદાં જુદાં સાધનો મુકાય છે જેમની હમણાં જરૂર નથી.

%%K t06-c3-03-1 p
3. --- આમ આપણે ઉપલા પાટિયા પર જોઈએ છીએ, \VUgras{શેકેલું ચિકન}
\textsuperscript{(7)}, \VUgras{માછલી} \textsuperscript{(8)} અને
\VUgras{મોટો પ્યાલો} \textsuperscript{(10)} \VUgras{ફળ}
\textsuperscript{(9)} સાથે. નીચે આ રહ્યું \VUgras{પનીર}
\textsuperscript{(11)}, \VUgras{બ્રેડ} \textsuperscript{(12)},
\VUgras{શીશીદાની} \textsuperscript{(13)} જેમાં કચુંબરને સ્વાદિષ્ટ
બનાવવા જોઈતું બધું છે : તેલ, સરકો, \VUgras{મીઠું}
\textsuperscript{(14)} અને \VUgras{મરી} \textsuperscript{(15)}. છેલ્લે,
નીચલા પાટિયા પર \VUgras{કેક} \textsuperscript{(17)} છે \VUgras{રાઈની
દાની} \textsuperscript{(16)} પાસે.

%%K t06-c3-04-1 p
4. --- ભોજનખંડનું ઘડિયાળ, જેને \VUgras{ભીંતઘડિયાળ}
\textsuperscript{(20)} કહે છે, તેનો આકાર ઘણો ખાસ છે. તે લાકડાના
ચોકઠામાં જડેલું છે જેમાં \VUgras{હવામાપક} \textsuperscript{(19)} અને
\VUgras{તાપમાપક} \textsuperscript{(18)} પણ છે.

%%K t06-tit-6 sub
\VUsaut{4.6mm}
\VUpk{10.2pt}{40}{રાતના ભોજનની પીરસણી.}
\VUsaut{3.0mm}

%%K t06-c3-05-1 p
5. --- ઓરડાની બધી ભીંતે \VUgras{લાકડાનાં પાટિયાં}
\textsuperscript{(21)} જડ્યાં છે, જે મીણ ચોપડેલા ઓકનાં છે. કાપડનો
\VUgras{બારણાનો પડદો} \textsuperscript{(22)} ઓસરીમાં જતું બારણું બરાબર
બંધ કરે છે. ત્યાં કુટુંબ રાતના ભોજન પછી કૉફી પીવા જશે. તુરત જ
\VUgras{પ્યાલા} \textsuperscript{(24)} અને \VUgras{રકાબીઓ}
\textsuperscript{(25)} \VUgras{વાસણના મેજ} \textsuperscript{(23)} પર
તૈયાર છે.

%%K t06-c3-06-1 p
6. --- મેજની ઉપર છતે \VUgras{લટકતો દીવો} \textsuperscript{(26)} જડ્યો
છે જેનો ઘણો તેજસ્વી દીવો સાંજે આખા ભોજનખંડને પ્રકાશિત કરે છે.

%%K t06-c3-07-1 p
7. --- હવે આપણે \VUgras{ભોજનમેજ} \textsuperscript{(27)} પોતે તપાસીએ.
કુટુંબ આવ્યું ત્યારે ભોજનનાં સાધનો તૈયાર હતાં. \VUgras{મેજપોશ}
\textsuperscript{(28)} પર, દરેક જમનારની જગ્યાએ, \VUgras{થાળી}
\textsuperscript{(33)}, \VUgras{નેપકિન} \textsuperscript{(29)},
\VUgras{ચમચો} \textsuperscript{(34)}, \VUgras{કાંટો}
\textsuperscript{(35)}, \VUgras{છરી} \textsuperscript{(36)} અને
\VUgras{પ્યાલો} \textsuperscript{(32)} મૂક્યાં હતાં. દ્રાક્ષાસવ માટે
\VUgras{બાટલીઓ} \textsuperscript{(30)} અને પાણી માટે \VUgras{પાણીની
બાટલી} \textsuperscript{(31)} પણ હતી.

%%K t06-c3-08-1 p
8. --- \VUgras{પીરસનારો} \textsuperscript{(39)} પહેલાં \VUgras{સૂપનું
વાસણ} \textsuperscript{(37)} લાવ્યો, જેમાં હજી થોડો સૂપ વરાળ કાઢે છે.
હવે તે પહેલી \VUgras{વાનગી} \textsuperscript{(38)} પીરસે છે, માંસની
વાનગી. પછી તે જે આપણે ગણાવી તે લાવશે, છેલ્લે, \VUgras{મીઠાઈની વાનગી}
પછી, પોતાના શેઠ ઓસરીમાં ગયા હશે તેનો લાભ લઈને મેજનાં સાધનો ઉપાડી લેશે
અને ભોજનખંડ ફરી ગોઠવશે.

%%K t06-c3-08-2 p
\textit{નોંધ.} --- \textit{ખોરાકને લગતા સામાન્ય શબ્દો અહીં ઘણા
ટૂંકમાં આપ્યા છે. યાદી બે કોષ્ટકો પર પૂરી કરવામાં આવશે, « હોટેલ »
(નં. ૧૩) અને « બજાર » (નં. ૧૫).}

%%K t06-tit-7 sub
\VUsaut{4.6mm}
\VUpk{10.2pt}{40}{રસોડું.}
\VUsaut{3.0mm}

%%K t06-c4-01-1 p
1. --- અડધા ખુલ્લા બારણામાંથી આપણે જે રસોડું જોવા લાગીએ છીએ, તેમાં
આપણે મનોહર અવ્યવસ્થા જોઈએ છીએ. \VUgras{ચૂલા} \textsuperscript{(1)}
આગળ, જેમાં \VUgras{આગ} \textsuperscript{(3)} ભડભડે છે, \VUgras{શેકવાનો
સળિયો} \textsuperscript{(5)} ગોઠવ્યો છે. સુંદર \VUgras{શેકેલું માંસ}
\textsuperscript{(7)} \VUgras{સળિયે પરોવેલું} \textsuperscript{(6)} છે
અને શેકાય છે.

%%K t06-c4-02-1 p
2. --- \VUgras{ફૂંકણી} \textsuperscript{(8)} અને \VUgras{કોલસાનો
પાવડો} \textsuperscript{(9)}, જે હમણાં વાપર્યાં, તે \VUgras{બળતણનાં
લાકડાંની} \textsuperscript{(10)} જેમ જમીન પર જ રહી ગયાં છે.

%%K t06-c4-03-1 p
3. --- સગડીની ઉપરનો ભાગ ગોઠવેલો છે ; \VUgras{ફાનસ}
\textsuperscript{(11)}, \VUgras{દીવાસળીની પેટી} \textsuperscript{(13)},
જેમાં \VUgras{દીવાસળી} \textsuperscript{(12)} છે, \VUgras{મીણબત્તીદાની}
\textsuperscript{(14)} અને \VUgras{નાની મીણબત્તીદાની}
\textsuperscript{(15)} \VUgras{મીણબત્તી} \textsuperscript{(16)} સાથે
પોતાની જગ્યાએ છે. પણ મોટી \VUgras{કોલસાની ડોલ}
\textsuperscript{(18)}, જે \VUgras{પથ્થરના કોલસાથી}
\textsuperscript{(17)} ભરેલી છે, જે \VUgras{રસોયણે}
\textsuperscript{(23)} હમણાં ભોંયરામાં ભરી, તે રસોડાની વચ્ચે જ રહી ગઈ
છે.

%%K t06-c4-04-1 p
4. --- ખરેખર, બિચારી સ્ત્રીએ ઉતાવળ કરવી પડે તેમ છે જેથી બપોરનું ભોજન
વખતસર તૈયાર થાય. હમણાં તે \VUgras{રસોઈના ચૂલા} \textsuperscript{(19)}
પાસે છે અને \VUgras{રસો} \textsuperscript{(24)} હલાવે છે, જે બહુ બળી ન
જવો જોઈએ.

%%K t06-c4-05-1 p
5. --- \VUgras{ભઠ્ઠીમાં} \textsuperscript{(20)} કેક શેકાય છે, જે તેણે
છોકરાંઓના નાસ્તા માટે તૈયાર કરી છે. રસોઈના ચૂલા ઉપર \VUgras{ડોયો}
\textsuperscript{(21)} અને \VUgras{ઝારો} \textsuperscript{(22)} લટકે છે.

%%K t06-tit-8 sub
\VUsaut{4.0mm}
\VUpk{10.2pt}{40}{વાસણ.}
\VUsaut{3.0mm}

%%K t06-c4-06-1 p
6. --- ઓરડાના છેડે લટકાવેલાં \VUgras{રસોઈનાં વાસણ}
\textsuperscript{(25)} ઘણાં સ્વચ્છ છે. સુંદર \VUgras{તાંબાની તપેલી}
\textsuperscript{(26)}, \VUgras{દૂધનું વાસણ} \textsuperscript{(27)},
\VUgras{ચાળણી} \textsuperscript{(28)} અને \VUgras{છીણી}
\textsuperscript{(29)} કાળજીથી સાફ કરેલાં છે.

%%K t06-c4-07-1 p
7. --- \VUgras{મીઠાની દાની} \textsuperscript{(30)} વાસણના કબાટ પર
\VUgras{ચાની કીટલી} \textsuperscript{(31)} અને \VUgras{તેલની મોટી
બાટલી} \textsuperscript{(32)} પાસે મૂકી છે. \VUgras{ખાનામાં}
\textsuperscript{(33)} કદાચ ભોજનનાં સાધનો અને રસોડાની છરીઓ છે.

%%K t06-c4-08-1 p
8. --- \VUgras{સાવરણી} \textsuperscript{(34)} અને \VUgras{પીંછાંનો
ઝૂડો} \textsuperscript{(35)} ખૂણામાં છે. રસોયણે હમણાં \VUgras{લાકડાનો
ટુકડો} \textsuperscript{(36)} વાપર્યો, તેણે તેને બારી પાસે મૂક્યો.

%%K t06-c4-09-1 p
9. --- \VUgras{હાથ લૂછવાનું} \textsuperscript{(37)} ભીંતે લટકે છે,
\VUgras{ગળણી} \textsuperscript{(38)} પાસે. રસોડાના આ ભાગમાં
\VUgras{શેકવાની જાળી} \textsuperscript{(39)}, \VUgras{માંસ કાપવાની છરી}
\textsuperscript{(40)} અને \VUgras{કાપવાનું પાટિયું}
\textsuperscript{(41)} મૂક્યાં છે. પછી, છરીની ઉપર,
\VUgras{છાજલી} \textsuperscript{(42)} પર, \VUgras{વાસણ}
\textsuperscript{(43)}, \VUgras{ઉકાળવાનું વાસણ} \textsuperscript{(44)}
પોતાના \VUgras{ઢાંકણ} \textsuperscript{(45)} સાથે અને \VUgras{મોટું
તપેલું} \textsuperscript{(46)} છે. \VUgras{કડાઈ} \textsuperscript{(47)}
પાસે લટકે છે.

%%K t06-c4-10-1 p
10. --- આ ખૂણામાં માત્ર તાંબાનો \VUgras{નળ} \textsuperscript{(48)} જ
ચમકે છે. \VUgras{વાસણ ધોવાની જગ્યા} \textsuperscript{(49)}, જેના પર
જાડો \VUgras{મોટો ઘડો} \textsuperscript{(50)} મૂક્યો છે, તે કદાચ
પોતાના નાના કબાટ સાથે ઘણી સગવડભરી છે, જેમાં \VUgras{સરકાનું પીપ}
\textsuperscript{(51)} પોતાના \VUgras{નળ} \textsuperscript{(52)} સાથે,
\VUgras{કચરાની પેટી} \textsuperscript{(53)} અને \VUgras{ગંદી થાળીઓ}
\textsuperscript{(54)} રખાય છે.

%%K t06-tit-9 sub
\VUsaut{4.0mm}
\VUpk{10.2pt}{40}{રસોડામાં મૂકેલી બીજી વસ્તુઓ.}
\VUsaut{3.0mm}

%%K t06-c4-11-1 p
11. --- \VUgras{મોટી કૂંડી} \textsuperscript{(55)}, જેમાં
\VUgras{શાકભાજી} \textsuperscript{(58)} ધોવાય છે, અને ઢાળેલા લોખંડનું
\VUgras{તપેલું} \textsuperscript{(56)}, જે \VUgras{લાદીના તળિયા}
\textsuperscript{(57)} પર મૂક્યું છે, પણ કદાચ એ કબાટમાં પોતાની જગ્યા
ધરાવે છે.

%%K t06-c4-12-1 p
12. --- રસોયણને ચૂલાની ખોટ નથી એ નક્કી, કેમ કે આ રહ્યો વળી, મેજના
ખૂણા પર, \VUgras{ગૅસનો ચૂલો} \textsuperscript{(59)} જેના પર નાની
તપેલીમાં પાણી ગરમ થાય છે. તેમાંથી નીકળતી \VUgras{વરાળ}
\textsuperscript{(60)} આપણને બતાવે છે કે તે કદાચ ઊકળે છે. કદાચ આ પાણી
કૉફી માટે વપરાશે, કેમ કે આ રહી મેજના બીજા છેડે, \VUgras{કૉફીની કીટલી}
\textsuperscript{(64)} અને \VUgras{કૉફીની ઘંટી} \textsuperscript{(63)}.

%%K t06-c4-13-1 p
13. --- ચૂલો \VUgras{ગૅસના બર્નર} \textsuperscript{(62)} સાથે લાંબી
\VUgras{રબરની નળી} \textsuperscript{(61)} વડે જોડ્યો છે.

%%K t06-c4-14-1 p
14. --- \VUgras{રોજની નોકરાણી} \textsuperscript{(66)}, જે રસોયણને મદદ
કરવા આવે છે, રાતના ભોજન માટે શાકભાજી તૈયાર કરે છે. તે સાફ અને ધોવાઈ
જશે ત્યારે તે તેમને \VUgras{કૂંડીમાં} \textsuperscript{(65)} મૂકશે અને
રાંધવાના સમયની રાહ જોશે.

%%K t06-tit-10 sub
\VUsaut{4.0mm}
{\centering\fontsize{10.2pt}{10.2pt}\selectfont\textls[40]{\scshape સ્નાનખંડ.} \textit{(નહાવાનો ઓરડો.)}\par}
\VUsaut{3.0mm}

%%K t06-c5-01-1 p
1. --- ઘરના મુખ્ય ઓરડા જાણવા માટે આપણે માત્ર એક જ અવિવેક કરવાનો બાકી
છે : સ્નાનખંડની ગોઠવણ જોવી. એમાં લાંબો સમય નહીં લાગે, કેમ કે તે મોટો
નથી, અને આપણે જલદી જાણી લઈશું કે \VUgras{નહાવાની ટબને}
\textsuperscript{(1)} સારો \VUgras{બંબો} \textsuperscript{(2)} છે, કે
\VUgras{સાબુદાની} \textsuperscript{(3)} નહાનારના હાથ પહોંચે એટલી પાસે
છે અને કે \VUgras{ટુવાલની દાંડી} \textsuperscript{(5)} લાગે છે કે
\VUgras{ટુવાલ} \textsuperscript{(4)} સહેલાઈથી સૂકવી શકશે.

%%K t06-c5-02-1 p
2. --- આ ઓરડાની ભીંતો \VUgras{ટાઇલ્સથી} \textsuperscript{(6)} ઢાંકેલી
છે, જેમણે \VUgras{ફુવારો} \textsuperscript{(7)} બેસાડવાનું શક્ય બનાવ્યું
છે, વાપરતી વખતે ઊડતું પાણી ભીંતોને બગાડશે એવો ડર રાખ્યા વિના. પગ
ધોવા માટે વપરાતી જસતની \VUgras{નાની કૂંડી} \textsuperscript{(8)}
રાચરચીલું પૂરું કરે છે.

\VUsaut{6.0mm}
\VUfilet{22mm}
